\clearpage\selectlanguage{english}\fancyhead[L]{\small POSTHUMOUS WRITINGS.}\pdfbookmark[0]{DETERMINATION OF OUR SERIES BY A DIFFERENTIAL EQUATION OF THE SECOND ORDER.}{nachlass216}
\SourceLeaf{216}
\begin{center}
\textbf{POSTHUMOUS WRITINGS.}\par\medskip
\textbf{DETERMINATION OF OUR SERIES}\par\medskip
\textbf{BY A DIFFERENTIAL EQUATION OF THE SECOND ORDER.}\par\medskip
\rule{0.08\linewidth}{0.3pt}\par\bigskip
\end{center}
\SixthArticle{38}
Setting \(F(\alpha,\beta,\gamma,x)=P\) for brevity, we have, by Article 4,
\[
\frac{\mathrm{d}P}{\mathrm{d}x}=\frac{\alpha\beta}{\gamma}
F(\alpha+1,\beta+1,\gamma+1,x)
\]
and hence, differentiating again,
\[
\frac{\mathrm{dd}P}{\mathrm{d}x^2}
=\frac{\alpha\beta(\alpha+1)(\beta+1)}{\gamma(\gamma+1)}
F(\alpha+2,\beta+2,\gamma+2,x)
\]
Equation IX of Article 10 therefore gives
\[
\text{[80]}\qquad
0=\alpha\beta P-(\gamma-(\alpha+\beta+1)x)
\frac{\mathrm{d}P}{\mathrm{d}x}-(x-xx)\frac{\mathrm{dd}P}{\mathrm{d}x^2}
\]
This differential equation of the second order may be regarded as a more precise definition of our function; but since \(P=F(\alpha,\beta,\gamma,x)\) is only a particular solution, rather than the general solution (as no constants have been introduced), the condition must be added that \(P\) begin with the value 1 for \(x=0\), and that, for the same value of \(x\), one also set
\(\frac{\mathrm{d}P}{\mathrm{d}x}=\frac{\alpha\beta}{\gamma}\)
and
\(\frac{\mathrm{dd}P}{\mathrm{d}x^2}
=\frac{\alpha\beta(\alpha+1)(\beta+1)}{\gamma(\gamma+1)}\).

Thus, for every value of \(x\) reached from the value \(x=0\) by continuous steps, without, however, reaching the value \(x=1\), for which \(x-xx=0\), \(P\) will be a completely determined quantity. But it is evident that \emph{in this way} one cannot reach positive real values of \(x\) greater than unity without passing through imaginary values.\SourceLeaf{217} Since this can be done in infinitely many different ways without violating continuity, it is not yet clear whether the same value of \(x\) may have several, or even infinitely many, distinct values of \(P\) corresponding to it, as is known to occur with several more familiar transcendental functions. We reserve a fuller discussion of this subject for later, since here we chiefly treat the case in which \(x\) is taken below, or at least not above, positive unity, and \(P\) is equal to the sum of the series \(F(\alpha,\beta,\gamma,x)\).

\SixthArticle{39}
In equation 80, writing \(1-y\) in place of \(x\) transforms it into
\[
\begin{aligned}&0=\alpha\beta P-(\alpha+\beta+1-\gamma-(\alpha+\beta+1)y)
\frac{\mathrm{d}P}{\mathrm{d}y}\\&-(y-yy)\frac{\mathrm{dd}P}{\mathrm{d}y^2}\end{aligned}
\]
which has a form similar to the former. Another particular solution follows at once:
\[
\begin{aligned}&P=F(\alpha,\beta,\alpha+\beta+1-\gamma,y)\\&=F(\alpha,\beta,\alpha+\beta+1-\gamma,1-x)\end{aligned}
\]
From this, by known principles, follows the general solution of equation 80,
\[
\begin{aligned}&\text{[81]}\qquad
P=MF(\alpha,\beta,\gamma,x)\\&+NF(\alpha,\beta,\alpha+\beta+1-\gamma,1-x)\end{aligned}
\]
where \(M,N\) denote arbitrary constants.

We also observe here in passing that it is easy to reduce to the form of equation 80 the following more general equation:
\[
0=AP+(B+Cy)\frac{\mathrm{d}P}{\mathrm{d}y}
+(D+Ey+Fyy)\frac{\mathrm{dd}P}{\mathrm{d}y^2}
\]
For let the roots of the equation \(0=D+Ey+Fyy\) be \(y=a,y=b\), or let \(D+Ey+Fyy\) be identically equal to the product \(F(y-a)(y-b)\). It is then evident that setting \(\frac{y-a}{b-a}=x\) and determining \(\alpha,\beta,\gamma\) so that
\[
\alpha\beta=\frac{A}{F},\quad
\alpha+\beta+1=\frac{C}{F},\quad
\gamma=-\frac{B+aC}{F(b-a)}
\]
transforms it into equation 80.

\SixthArticle{40}
With the aid of differential equation 80, several very remarkable theorems about our series can be derived, some general and some more special; and we do not doubt that many more, and more substantial ones, still lie hidden, reserved for further investigation. Let us present here those that we have so far been able to uncover.
\par\SourceLeaf{218}
Let us set $P=(1-x)^\mu P'$; then
\[
\begin{aligned}
\frac{\mathrm{d}P}{\mathrm{d}x}&=-\mu(1-x)^{\mu-1}P'+(1-x)^\mu\frac{\mathrm{d}P'}{\mathrm{d}x}\\
\frac{\mathrm{dd}P}{\mathrm{d}x^2}&=\mu(\mu-1)(1-x)^{\mu-2}P'-2\mu(1-x)^{\mu-1}\frac{\mathrm{d}P'}{\mathrm{d}x}\\&\qquad+(1-x)^\mu\frac{\mathrm{dd}P'}{\mathrm{d}x^2}
\end{aligned}
\]
Substituting these values in equation 80 and dividing by $(1-x)^{\mu-1}$ gives
\[
\begin{aligned}
0={}&P'\{\alpha\beta(1-x)+(\gamma-(\alpha+\beta+1)x)\mu-x(\mu\mu-\mu)\}\\
&-\frac{\mathrm{d}P'}{\mathrm{d}x}\{(\gamma-(\alpha+\beta+1)x)-2\mu x\}(1-x)\\&-\frac{\mathrm{dd}P'}{\mathrm{d}x^2}\{x-xx\}(1-x)
\end{aligned}
\]
Let us determine $\mu$ so that the multiplier of $P'$ is divisible by $1-x$. This is achieved by setting either $\mu=0$ or $\mu=\gamma-\alpha-\beta$. The first assumption would teach us nothing new, but substitution of the latter value gives
\[
\begin{aligned}&0=P'\{\alpha\beta-\alpha\gamma-\beta\gamma+\gamma\gamma\}-\frac{\mathrm{d}P'}{\mathrm{d}x}\{\gamma-(2\gamma-\alpha-\beta+1)x\}\\&-\frac{\mathrm{dd}P'}{\mathrm{d}x^2}\{x-xx\}\end{aligned}
\]
or
\[
\begin{aligned}&0=P'(\gamma-\alpha)(\gamma-\beta)-\frac{\mathrm{d}P'}{\mathrm{d}x}\{\gamma-((\gamma-\alpha)+(\gamma-\beta)+1)x\}\\&-\frac{\mathrm{dd}P'}{\mathrm{d}x^2}(x-xx)\end{aligned}
\]
which has exactly the same form as equation 80. Thus, since for $x=0$ we evidently have $P'=1$ and $\frac{\mathrm{d}P'}{\mathrm{d}x}=\frac{\alpha\beta}{\gamma}-\mu=\frac{(\gamma-\alpha)(\gamma-\beta)}{\gamma}$, it is clear that its solution is $P'=F(\gamma-\alpha,\gamma-\beta,\gamma,x)$, so that in general
\[
\text{[82]}\qquad F(\gamma-\alpha,\gamma-\beta,\gamma,x)=(1-x)^{\alpha+\beta-\gamma}F(\alpha,\beta,\gamma,x)
\]
This is the source of the transformation of the series
\[
1+\frac{2\,.\,8}{9}x+\frac{3\,.\,8\,.\,10}{9\,.\,11}xx+\frac{4\,.\,8\,.\,10\,.\,12}{9\,.\,11\,.\,13}x^3+\text{etc.}=F(2,4,\tfrac{9}{2},x)
\]
into
\[
\begin{aligned}&(1-x)^{-\frac{3}{2}}\left(1+\frac{1\,.\,5}{2\,.\,9}x+\frac{1\,.\,3\,.\,5\,.\,7}{2\,.\,4\,.\,9\,.\,11}xx+\text{etc.}\right)\\&=(1-x)^{-\frac{3}{2}}F(\tfrac{5}{2},\tfrac{1}{2},\tfrac{9}{2},x)\end{aligned}
\]
which we had stated without proof in the Berlin Astronomical Ephemerides for 1814, p. 257 [Supplement to Articles 90 and 100 of the Theoria motus].

